\clearpage
\setcounter{section}{5}

\section{Spezifische Wärme und Wärmeleitung}
\subsection{Zustandsdichte in 3D (isotroper Körper $\hat{=}$ BZ $\to $ Kugel)}
Wir betrachten $N_{i}^3$ primitive Elementarzellen in einem Würfel mit der Kantenlänge $L$. Das Volumen der Probe ist $L^3 = V$ und es gilt die periodische Randbedingung  
$$
k_{x},\, k_{y},\, k_{z} = 0, \, \pm \frac{2\pi}{L}, \, \pm \frac{4\pi}{L}, \, \ldots,\, \pm \frac{2\pi N_{i}}{L}.
$$

\begin{figure}[H]
    \centering
    \incfig{vl23_abbildung_1}
    \label{fig:vl23_abbildung_1}
\end{figure}

Der Platz für ein $k$ in 3D ist
$$
\left( \frac{2\pi}{L} \right) ^3 \quad \text{bzw.} \quad
\left( \frac{L}{2\pi} \right) ^3 = \frac{V}{8\pi^3}
$$ 
erlaubte $k$-Werte pro Einheitsvolumen im $\Vec{k}$-Raum für jede Polarisation und jeden Zweig der Dispersionsrelation.\\
Wir betrachten eine Kugel mit Radius $k$. Die Gesamtzahl der Zustände mit einem Wellenvektor dessen Betrag kleiner ist als $k$ ergibt sich durch
$$
N = \frac{\frac{4\pi k^3}{3}\, \text{(Volumen der Kugel)}}{\left( \frac{2\pi}{L} \right) ^3 \, \text{(Volumen für einen k-Zustand)}} = \left( \frac{L}{2\pi}^3 \right) \frac{4\pi k^3}{3} = \frac{Vk^3}{6\pi^2} = \frac{Vk^3(\omega)}{6\pi^2}
$$

Die Zustandsdichte ergibt sich daraus mit
\begin{equation}
	\label{eq:II.6.1}
	D(\omega) = \frac{\mathrm{d}N}{\mathrm{d}\omega} = \frac{Vk^2}{2\pi^2} \frac{\mathrm{d} k}{\mathrm{d}\omega}.
\end{equation}

\subsection{Allgemeine Behandlung von $D(\omega)$}
Die Anzahl der erlaubten $k$-Werte mit Phononenfrequenzen zwischen $\omega $ und $\omega + \mathrm{d} \omega$ ist
$$
D(\omega) \text{d}\omega = \left( \frac{L}{2\pi} \right) ^3 \int_{\omega=\text{const.}}^{\omega+\mathrm{d}\omega=\text{const}} \hspace{-50pt}\mathrm{d}^{3}k,
$$ 
wobei sich das Integral über eine Schale erstreckt, die von zwei Oberflächen im $k$-Raum eingeschlossen wird, auf denen die Phononenfrequenz konstant ist ($\omega$ und $\omega + \mathrm{d}\omega$)

\begin{figure}[H]
    \centering
    \def\svgwidth{.499\columnwidth}
    \subimport{figures/}{vl23_abbildung_2.pdf_tex}
    \label{fig:vl23_abbildung_2}
\end{figure}

Das Volumenelement zwischen den Schalen ist ein Zylinder mit $\mathrm{d}^3k = \mathrm{d}S_{\omega} \cdot \mathrm{d} k_{\perp}$, es folgt
$$
\int_{\text{Schale}} \mathrm{d}^3 k = \int \mathrm{d}S_{\omega} \mathrm{d}k_{\perp}.
$$
Mit 
$$
\left| \mathrm{grad}_{k}(\omega) \right| = \left| \nabla _{k} (\omega) \right| = \frac{\mathrm{d}\omega}{\mathrm{d} k_{\perp}}
$$ 
folgt
$$
\mathrm{d} S_{\omega} \cdot \mathrm{d}k_{\perp} = \mathrm{d} S_{\omega} \cdot \frac{\mathrm{d}\omega}{\left| \nabla _{k} (\omega) \right| } = \mathrm{d}S_{\omega} \cdot \frac{\mathrm{d}\omega}{v_\mathrm{g}}
$$ 
wobei $v_{\mathrm{g}} = \left| \nabla _{k} (\omega) \right| $ der Betrag der Gruppengeschwindigkeit des Phonons ist. Für die Zustandsdichte ergibt sich der Ausdruck
$$
D(\omega) \mathrm{d}\omega = \left( \frac{L}{2\pi} \right) ^3 \int \frac{\mathrm{d}S_{\omega}}{v_{\mathrm{g}}} \mathrm{d} \omega.
$$ 

\begin{equation}
	\label{eq:II.6.2}
	D(\omega) = \left( \frac{L}{2\pi} \right) ^3 \int \frac{\mathrm{d}S_{\omega}}{v_{\mathrm{g}}}
\end{equation}

Diese Gleichung ermöglicht die Bestimmung der Zustandsdichte aus der Dispersionsrelation und macht erkennbar, dass sich für flache Stellen der Dispersionsrelation hohe Zustandsdichten ergeben, da $\left| \nabla _{k} (\omega) \right| $ klein wird. Es bestehen sogenannte Van-Hove Singularitäten für $\left| \nabla _{k}(\omega) \right| \to 0$.\\

\paragraph{Anmerkung}\,\\ 
Bei 1D divergiert die Zustandsdichte. Bei 3D hingegen divergiert der Integrand, aber nicht das Integral \ref{eq:II.6.2}, es ergeben sich Spitzen in der Zustandsdichte.\\

Ein Nachteil von Gleichung \ref{eq:II.6.2} ist, dass die Disperionsrelation $\omega(k)$ bekannst sein muss, wobei diese sehr stark systemabhängig ist. Wünschenswert wäre es also aus \ref{eq:II.6.2} allgemeingültige Informationen bzw. Erkenntnisse zu ziehen. Dies ist möglich für niedrige ($T \le \SI{300}{\kelvin}$) Temperaturen. 

\paragraph{Debye'sche Modell der Zustandsdichte}\,\\ 
Der Kristall befindet sich im thermischen Gleichgewicht, d.h. Phonenen werden nur thermisch angeregt also $k_{\mathrm{B}}T \le \omega_{\text{opt}} (k)$. Es werden also nur akustische Schwingungsmoden im Zentrum der BZ angeregt.\\
Dort gilt: $\omega(k) = c_{i} k$ mit $i$ = Zweig der Dispersion (1 longitudinal, 2 transversal) und der Schallgeschwindigkeit $c_{i}$. \\
Daraus folgt
\begin{enumerate}
	\item $\left| \nabla _{\Vec{k}} (\omega) \right| = c_{i}$ 
	\item $\omega$ = konst. ist kugelfläche im $k$-Raum mit Radius
		 \begin{equation}
		 	\label{eq:II.6.3}
		 	k = \omega / c_{i}
		 \end{equation} 
\end{enumerate}
Setzt man Gleichung \ref{eq:II.6.1} oder \ref{eq:II.6.3} in \ref{eq:II.6.2} ein ergibt sich
$$
D_{i}(\omega) \mathrm{d}\omega = \frac{V}{\left( 2\pi \right) ^3} \mathrm{d} \omega \frac{1}{c_{i}} 4 \pi k^2
$$ 
\begin{equation}
	\label{eq:II.6.4}
	D_{i} (\omega) \mathrm{d}\omega = \frac{V}{2\pi^2} \frac{\omega^2}{c_{i}^3} \mathrm{d}\omega
\end{equation}

Insgesamt ergibt sich für die Zustandsdichte
$$
D(\omega) = \frac{V}{2\pi^2} \omega^2 \left( \frac{1}{c_{\mathrm{L}}^3} + \frac{1}{c_{\mathrm{T_1}}^3} + \frac{1}{c_{\mathrm{T_2}}^3} \right).
$$ 
Mit der Näherung $c_{\mathrm{L}} = c_{\mathrm{T1}}= c_{\mathrm{T2}} = v_{\mathrm{g}}$ folgt

\begin{equation}
	\label{eq:II.6.5}
	D(\omega) = \frac{3v \omega^2}{2 \pi^2 v_{^{g}}^3}
\end{equation}

\begin{figure}[H]
    \centering
    \def\svgwidth{.41\columnwidth}
    \subimport{figures/}{vl23_abbildung_3.pdf_tex}
    \label{fig:vl23_abbildung_3}
\end{figure}

Problem:
$$
\int_{0}^{\infty} D(\omega) \mathrm{d}\omega = \infty
$$ 
Lösung: $D(\omega)$ abschneiden bei der \textbf{Debye-Frequenz} $\mathbf{\omega_{D}}$. Es folgt

\begin{equation}
	\label{eq:II.6.6}
	\int_{0}^{\omega_{D}} D(\omega) \mathrm{d}\omega = 3N
\end{equation} 
was der Zahl der Freiheitsgrade (Schwingungsmoden) entspricht mit der Zahl der Einheitszellen $N$ bei einatomiger Basis. Die Debye-Frequenz berechnet sich mittels
\begin{equation}
	\label{eq:II.6.7}
	\omega_{D} = v_{\mathrm{g}} \sqrt[3]{ \frac{6\pi^2N}{V}} 
\end{equation}
und die Debye-Temperatur $\Theta_{D}$ 
\begin{equation}
	\label{eq:II.6.8}
	\hbar \omega_{D} = k_{\mathrm{B}}\Theta_{D}.
\end{equation}

\subsection{Spezifische Wärmekapazität des Kristallgitters}
\paragraph{Klassische Behandlung}\,\\ 
\begin{itemize}
	\item Die Molwärme, also die Wärme welche nötig ist um \SI{1}{\mol} um \SI{1}{\kelvin} zu erwärmen, berechnet sich mittels
		\begin{equation}
			\label{eq:II.6.9}
			C_{V} = \left( \frac{\mathrm{d}Q}{\mathrm{d}T} \right) _{V} = \frac{\mathrm{d}U}{\mathrm{d}T}
		\end{equation}
		mit der inneren Energie $U$. $C_{V}$ wird berechnet, $C_{p}$ wird gemessen. Im Festkörper gilt $C_{V} \approx C_{p}$ auf eine Genauigkeit von $1\,\text{\textperthousand}$.
	
	\item Nach dem Gleichverteilungssatz der klassischen statistischen Mechanik entspricht $k_{\mathrm{B}}T$ der mittleren Energie pro Schwingungsfreiheitsgrad ($E_{\text{kin}} + E_{\text{pot}}$) ($\left( 2 \cdot \frac{1}{2} k_{\mathrm{B}}T =  \right) $)
	
	\item Im Festkörper gilt bei $N$ Atomen mit 3 Schwingungsfreiheitsgraden pro Atom $U\left( V,T \right) = 3N k_{\mathrm{B}}T$, und somit die Dulong-Petit'shce Gesetz
		\begin{equation}
			\label{eq:II.6.10}
			C_{V} = \left( \frac{\partial U}{\partial T} \right) _{V} = 3Nk_{\mathrm{B}} 
		\end{equation}
		molare Wärmekapazität $N = N_{A}$, $N_{A}k_{\mathrm{B}} = R$ 
		$$
		\to C_{V} = 3R = \SI{24,94}{\frac{\joule}{\mol \cdot \kelvin}} = \text{konst.}
		$$ 
		Im Experiment werden bei tiefen Temperaturen starke Abweichungen beobachtet!\\
		Bei Isolatoren: $C_{V}\sim T^3$, d.h. $T \to 0$ folgt $C_{V} \to 0$ (Phononen)\\
		Bei Metallen: $C_{V} \sim AT + BT^3$, d.h. $T\to 0$ folgt $C_{V} \to 0$ (Phonon Metallelektronen)
\end{itemize}

\paragraph{QM Behauptung}\,\\
Ausdruck fúr die innere Energie $U$ 
\begin{equation}
	\label{eq:II.6.11}
	U = \int \mathrm{d} \omega D(\omega) \varepsilon(\omega,T)
\end{equation}
\begin{itemize}
	\item $\varepsilon(\omega,T)$ = mittlere Energie der Schwingungsmode $\omega$ 
	\item Energie des quantenmechanischen Oszillators 
		$$
		E_{n} = \hbar \omega\left( n + \frac{1}{2} \right) 
		$$ 
	\item Erwartungswert der Besetzungszahl $\left<n \right>$ 
		\begin{equation}
			\label{eq:II.6.12}
			\varepsilon(\omega,T) = \hbar \omega\left( \left<n \right> + \frac{1}{2} \right) 
		\end{equation}
		\begin{equation}
			\label{}
			\varepsilon(\omega,T) = \hbar \omega \left( \frac{1}{2} + \frac{1}{e^{\hbar \omega / (k_{\mathrm{B}} T)}-1} \right) 
		\end{equation}
		Bose-Statistik
		\begin{equation}
			\label{}
			\left<n \right>_{\text{therm}} = \frac{1}{e^{\hbar \omega / (k_{\mathrm{B}}T) } - 1}
		\end{equation}
		Der Erwartungswert der Quantenzahl $n$ für einen Oszillator im thermischen Gleichgewicht enspricht der mittleren Phononenzahl im thermischen Gleichgewicht.
\end{itemize}
